\section{Theoretical Contributions}
\label{sec:contributions}

This section derives the two main theorems and two propositions that characterise the
performance guarantees of \sys{}.
We additionally state two corollaries that follow immediately from the main theorems,
and provide extended discussion of the design implications of each result.

\subsection{Theorem 1: Repair Convergence of the \tvg{} Cycle}

The first theorem characterises the probability of successfully compiling a TikZ figure
within a bounded number of LLM repair attempts.
This is the foundational reliability guarantee of the Typed Visual Grammar system.

\begin{theorem}[\tvg{} Repair Convergence]
\label{thm:repair}
Let $p\in(0,1)$ be the probability that a single LLM repair attempt produces a
TikZ figure that compiles without error under $L$ (the \LaTeX{} compiler oracle),
and let $K\geq 0$ be the maximum number of \emph{additional} repair attempts
permitted after the initial generation attempt.
Then the probability of at least one successful compilation within $K+1$ total attempts is
\[
  P(\mathrm{success}\mid K) \;=\; 1 - (1-p)^{K+1}.
\]
\end{theorem}

\begin{proof}
Each of the $K+1$ attempts is modelled as an independent Bernoulli trial with success
probability $p$.
Independence holds because each repair attempt involves a fresh LLM call in which
the complete TikZ source code is regenerated from the blueprint template, conditioned
on the most recent stderr message from $L$.
Since the LLM does not maintain internal state between API calls (the session context
is supplied anew in each call), the outcomes of successive attempts are independent
random variables $X_1,\ldots,X_{K+1}\stackrel{\mathrm{iid}}{\sim}\mathrm{Bernoulli}(p)$.

The event $\{\mathrm{success}\}$ is the complement of the event $\{\text{all attempts
fail}\} = \{\forall\, i: X_i=0\}$.
By independence:
\[
  P(\text{all fail}) = \prod_{i=1}^{K+1}P(X_i=0) = (1-p)^{K+1}.
\]
Therefore, $P(\mathrm{success}\mid K) = 1 - (1-p)^{K+1}$. \qed
\end{proof}

\begin{corollary}[Practical Convergence at Default Settings]
\label{cor:repair}
For $p=0.72$ (empirically measured on our figure corpus, Section~\ref{sec:experimental})
and $K=3$ (the configured attempt budget):
\[
  P(\mathrm{success}\mid 3) = 1 - (1-0.72)^4 = 1 - 0.28^4 = 1 - 0.00615 \approx 0.9938.
\]
With probability greater than 99.3\%, the \tvg{} cycle produces a compilable figure
within four attempts.
\end{corollary}

The theorem makes the trade-off between attempt budget and success probability precise
and actionable for a system designer.
A designer who requires $P(\mathrm{success}\mid K)\geq 1-\varepsilon$ for a target
reliability $\varepsilon>0$ must choose:
\[
  K \;\geq\; \left\lceil \frac{\log\varepsilon}{\log(1-p)} \right\rceil - 1.
\]
For $\varepsilon=10^{-3}$ (one failure per thousand figures) and $p=0.72$, this gives
$K\geq 7$ (eight total attempts).
For $\varepsilon=10^{-4}$ (one per ten thousand), $K\geq 10$.
The current default $K=3$ delivers $\varepsilon\approx 6.2\times 10^{-3}$, which
is acceptable for interactive academic document generation where the user can be
prompted to manually supply a figure for the rare failure case.
For unattended batch processing where manual intervention is not possible, $K$ should
be increased to at least 7.

A second implication of Theorem~\ref{thm:repair} concerns the \emph{sensitivity} of
$P(\mathrm{success}\mid K)$ to $p$.
At $K=3$, the gradient $\partial P/\partial p = 4(1-p)^3$ evaluates to
$4\times 0.28^3 \approx 0.088$ at $p=0.72$.
A 10-percentage-point improvement in per-attempt success probability (from 0.72 to 0.82)
would increase $P(\mathrm{success}\mid 3)$ from $99.4\%$ to $99.9\%$, a gain of
only $0.5$ percentage points.
This reveals a \emph{diminishing returns} structure: once per-attempt success exceeds
approximately $0.70$, additional improvements in $p$ yield minimal gains at $K=3$.
The more effective lever for high-reliability deployments is increasing $K$.

\subsection{Theorem 2: Order Preservation and Amdahl Speedup}

The second theorem establishes two properties of the \texttt{concurrentMap} primitive
that together underpin Invariant~I4 and the scalability claim.

\begin{theorem}[Order-Preserving \texttt{concurrentMap} and Amdahl Speedup]
\label{thm:concurrent}
Let $\mathcal{F}=\{f_1,\ldots,f_n\}$ be a list of independent document-section
generation tasks, and let $\texttt{concurrentMap}(\mathcal{F},N)$ execute them
with at most $N$ workers operating in parallel.
Let $\alpha\in(0,1)$ be the fraction of total wall-clock time that is inherently
sequential (single-worker: billing accumulation, index-order assembly).
Then:
\begin{enumerate}[label=(\alph*)]
  \item \textbf{Order preservation.} The output list emits $f_1(\cdot),\ldots,f_n(\cdot)$
        in index order, regardless of the order in which workers complete.
  \item \textbf{Amdahl speedup.} The wall-clock speedup over sequential execution is
        \[
          S(N) \;=\; \frac{1}{\alpha + (1-\alpha)/N}.
        \]
\end{enumerate}
\end{theorem}

\begin{proof}
\textbf{(a) Order preservation.}
The \texttt{concurrentMap} primitive allocates a pre-sized output buffer
$B[1\ldots n]$ indexed by task position before any worker is started.
Worker allocation is done by a coordinator that maintains a task queue
$Q=\langle f_1,f_2,\ldots,f_n\rangle$ from which workers pull tasks in index order.
When worker $w$ completes task $f_i$, it writes its result to $B[i]$.
Because tasks are independent (no inter-task data dependencies by assumption), and
because each $B[i]$ is written exactly once by the unique worker assigned to $f_i$,
no write conflict can occur.

After all $n$ tasks have signalled completion via a counting semaphore
(decremented from $n$ to $0$), the coordinator reads $B$ sequentially from index
$1$ to $n$ and emits the results.
The emission order is therefore exactly the index order $B[1],B[2],\ldots,B[n]$,
regardless of the order in which workers completed.
Crucially, the coordinator does \emph{not} emit $B[j]$ until all $B[1\ldots j-1]$
have been emitted, even if $B[j]$ is filled before $B[j-1]$.
The semaphore ensures no premature emission. \qed

\textbf{(b) Amdahl speedup.}
Let $T_{\mathrm{seq}}$ be the total wall-clock time for sequential execution.
Of this, $\alpha T_{\mathrm{seq}}$ is inherently sequential and $(1-\alpha)T_{\mathrm{seq}}$
is perfectly parallelisable.
With $N$ workers, the parallel fraction takes time $(1-\alpha)T_{\mathrm{seq}}/N$,
and the sequential fraction remains $\alpha T_{\mathrm{seq}}$.
Total time: $\alpha T_{\mathrm{seq}} + (1-\alpha)T_{\mathrm{seq}}/N
            = T_{\mathrm{seq}}(\alpha + (1-\alpha)/N)$.
Speedup: $S(N) = T_{\mathrm{seq}} / [T_{\mathrm{seq}}(\alpha+(1-\alpha)/N)]
              = 1/(\alpha+(1-\alpha)/N)$. \qed
\end{proof}

\begin{corollary}[Stage-Specific Speedup]
\label{cor:amdahl}
For Stage~2 (reference fetch, dominated by LLM-wait, measured $\hat{\alpha}=0.15$)
with $N=4$ workers:
\[
  S(4) = \frac{1}{0.15 + 0.85/4} = \frac{1}{0.3625} \approx 2.76.
\]
Four concurrent reference-fetch workers achieve 2.76$\times$ wall-clock speedup.
The theoretical maximum (as $N\to\infty$) is $1/\alpha = 1/0.15 \approx 6.67$.
At $N=4$, the system recovers $\frac{S(4)-1}{S(\infty)-1} = \frac{1.76}{5.67} \approx 31\%$
of the available parallelism headroom.
\end{corollary}

Part~(a) of Theorem~\ref{thm:concurrent} is non-trivial precisely because of
heterogeneous task durations.
In a system where all tasks take identical time, any ordering strategy produces correct
index order.
But in the \sys{} pipeline, section drafting tasks have highly variable durations:
a section with many references to reconcile may take three times longer than a short
methodological section.
Without the buffer-and-semaphore mechanism, a fast worker completing a later section
would emit its result before a slow worker completes an earlier section, violating
Invariant~I4.
The proof of part~(a) establishes that the buffer-and-semaphore mechanism is exactly
the correct engineering pattern for this scenario.

Figure~\ref{fig:amdahl} shows the Amdahl speedup curves for four values of $\alpha$,
with three annotated operating points corresponding to the three parallel stages in
the platform (Stage~2 at $N=4$, Stage~2.5 at $N=3$, and Stage~3 at $N=2$).

\begin{figure}[h!]
\centering
\begin{tikzpicture}
\begin{axis}[PL,
  xlabel={Concurrent workers $N$},
  ylabel={Speedup $S(N)$},
  title={Amdahl's Law applied to \texttt{concurrentMap} parallel stages},
  xmin=1, xmax=16.5, ymin=1, ymax=9.5,
  xtick={1,2,4,6,8,10,12,14,16},
  legend pos=north west,
  width=13cm, height=7cm,
]
\addplot[C4, samples=60, domain=1:16]{1/(0.50+0.50/x)};
\addlegendentry{$\alpha=0.50$ (I/O bound)}
\addplot[C2, samples=60, domain=1:16]{1/(0.30+0.70/x)};
\addlegendentry{$\alpha=0.30$}
\addplot[C3, samples=60, domain=1:16]{1/(0.15+0.85/x)};
\addlegendentry{$\alpha=0.15$ (Stage\,2, LLM-wait)}
\addplot[C1, samples=60, domain=1:16]{1/(0.05+0.95/x)};
\addlegendentry{$\alpha=0.05$ (ideal)}
% Fill between ideal and sequential bounds
\addplot[name path=ID, draw=none, samples=60, domain=1:16]{1/(0.05+0.95/x)};
\addplot[name path=SQ, draw=none, samples=60, domain=1:16]{1/(0.50+0.50/x)};
\addplot[C6, opacity=0.18] fill between[of=ID and SQ];
% Operating points
\addplot[C3, mark=*, only marks, mark size=5pt] coordinates{(4,{1/(0.15+0.85/4)})};
\node[font=\tiny\bfseries, text=C3, anchor=west]
  at (axis cs:4.3,{1/(0.15+0.85/4)}){Stage\,2 $N\!=\!4$};
\addplot[C5, mark=square*, only marks, mark size=4.5pt]
  coordinates{(3,{1/(0.30+0.70/3)})};
\node[font=\tiny\bfseries, text=C5, anchor=west]
  at (axis cs:3.3,{1/(0.30+0.70/3)}){Stage\,2.5 $N\!=\!3$};
\addplot[C2, mark=triangle*, only marks, mark size=4.5pt]
  coordinates{(2,{1/(0.50+0.50/2)})};
\node[font=\tiny\bfseries, text=C2, anchor=east]
  at (axis cs:1.7,{1/(0.50+0.50/2)}){Stage\,2.3 $N\!=\!2$};
\end{axis}
\end{tikzpicture}
\caption{Amdahl's Law speedup curves for the order-preserving
         \texttt{concurrentMap} primitive (Theorem~\ref{thm:concurrent}).
         Shaded region spans achievable speedup between $\alpha\!=\!0.05$
         and $\alpha\!=\!0.50$. The three annotated operating points correspond
         to the configured concurrency bounds of each parallel pipeline stage.}
\label{fig:amdahl}
\end{figure}


\subsection{Proposition 1: Token Continuation Expectation}

\begin{proposition}[Token Continuation Expectation]
\label{prop:token}
Let $W_K$ be the total tokens consumed across up to $K+1$ repair attempts in the \tvg{}
cycle, where the cycle terminates at the first successful attempt.
Assume each attempt consumes a fixed token budget $w$ regardless of success or failure.
Then:
\[
  \mathbb{E}[W_K] \;=\; w\sum_{k=0}^{K}(1-p)^k \;=\; W_\infty\bigl(1 - e^{-\lambda K}\bigr)
\]
where $W_\infty = w/p$ is the expected total spend under unlimited retries,
and $\lambda = -\log(1-p)$ is the per-attempt failure rate.
\end{proposition}

\begin{proof}
Let $N_K$ be the number of attempts made before the first success, truncated at $K+1$.
The probability that exactly $k+1$ attempts are made (the first $k$ fail, the $(k+1)$-th
succeeds, where $k\leq K-1$) is $p(1-p)^k$.
The probability that all $K+1$ attempts are made (all fail, or the last succeeds) is
$(1-p)^K$.

\begin{align*}
\mathbb{E}[W_K] &= w\,\mathbb{E}[N_K] \\
&= w\left[\sum_{k=0}^{K-1}(k+1)p(1-p)^k + (K+1)(1-p)^K\right]
\end{align*}

This simplifies by observing that regardless of when the cycle terminates,
every attempt up to termination consumes exactly $w$ tokens.
An equivalent computation counts the expected number of attempts:
\[
\mathbb{E}[N_K] = \sum_{k=0}^{K}\Pr[\text{at least }k+1\text{ attempts made}]
= \sum_{k=0}^{K}(1-p)^k = \frac{1-(1-p)^{K+1}}{p}.
\]
Multiplying by $w$:
\[
\mathbb{E}[W_K] = \frac{w(1-(1-p)^{K+1})}{p} = W_\infty(1-(1-p)^{K+1}).
\]
Using $(1-p) = e^{-\lambda}$, we have $(1-p)^{K+1} \approx e^{-\lambda K}$ for large $K$,
yielding the stated expression. \qed
\end{proof}

Proposition~\ref{prop:token} has a direct implication for the Billing Operator
coefficient.
For $p=0.72$ and $w=266\,\text{PMU}$:
$W_\infty = 266/0.72 \approx 369\,\text{PMU}$.
At $K=3$:
$\mathbb{E}[W_3] = 369\times(1-0.28^4) \approx 369\times 0.9938 \approx 367\,\text{PMU}$.
The per-figure billing coefficient of $800\,\text{PMU}$ in $\mathcal{B}$ is
$\approx 2.2\times W_\infty$, providing a safety margin of $2.2\times$
above the expected token expenditure.
This margin accommodates the stochastic variance in token consumption across figure
types and ensures that the billing accumulation event at figure completion always
records a value $\leq 800\,\text{PMU}$ with high probability.

\subsection{Proposition 2: Pipeline as a Markov Decision Process}

\begin{proposition}[Pipeline as MDP]
\label{prop:mdp}
The pipeline $\Pi$ can be modelled as a finite Markov Decision Process
$\mathrm{MDP}=(\mathcal{S},\mathcal{A},\mathcal{T},\mathcal{R},\gamma)$ where:
\begin{itemize}
  \item $\mathcal{S}=\{S_0,S_1,S_2,S_2',S_3,S_4,S_5,S_{\mathrm{fail}}\}$
        is the state space (document states plus an absorbing failure state);
  \item $\mathcal{A}=\{a_1,a_2,a_{2.5},a_3,a_4,a_5\}$ is the action set
        (stage-transition decisions);
  \item $\mathcal{T}(S_i\mid S_{i-1},a_i)=1-\delta_i$ (successful transition)
        and $\mathcal{T}(S_{\mathrm{fail}}\mid S_{i-1},a_i)=\delta_i$
        (failure transition), where $\delta_i$ is the per-stage failure probability;
  \item $\mathcal{R}(S_{i-1},a_i)=-\mathcal{B}_i$ is the stage-local reward,
        where $\mathcal{B}_i$ is the marginal billing cost of stage $i$;
  \item $\gamma=1$ (no time discounting, since session duration is bounded).
\end{itemize}
Under the deterministic policy $\pi^*$ that always executes the next stage in sequence,
the expected cumulative reward is $-\sum_{i=1}^{5}\mathcal{B}_i = -\mathcal{B}(T,n_{\mathrm{img}})$.
\end{proposition}

\begin{proof}[Proof sketch]
The Markov property holds because each stage's output is a deterministic function
of its input state: $S_i = F_i(S_{i-1})$ conditioned on success, or $S_{\mathrm{fail}}$
conditioned on failure.
There is no hidden state beyond the current document state.
The reward function $\mathcal{R}(S_{i-1},a_i)=-\mathcal{B}_i$ captures the cost of
executing stage $i$, negated so that cost minimisation corresponds to reward maximisation.
By additivity of $\mathcal{B}$, the sum of stage-local rewards equals the total billing
cost, negated. \qed
\end{proof}

Proposition~\ref{prop:mdp} is primarily of theoretical interest as the foundation
for future learned-policy extensions.
In the current system, $\pi^*$ is fixed: always execute the next stage.
This is optimal when all stage-local failure probabilities $\delta_i$ are small
(which holds in the evaluated corpus, see Table~\ref{tab:invariants}).
When $\delta_i$ is large — for example, when the reference-fetch endpoint is
intermittently unavailable — a learned policy could adaptively skip $F_2$ and
proceed with a reduced reference set rather than aborting the session entirely.
Proposition~\ref{prop:mdp} formalises the decision problem that such a policy would need to solve.

\subsection{Summary of Formal Guarantees}

Table~\ref{tab:formal_guarantees} summarises the formal guarantees established
by the definitions, theorems, and propositions in this section.

\begin{table}[h!]
\centering
\caption{Summary of formal guarantees for the \sys{} system $\mathcal{P}$.}
\label{tab:formal_guarantees}
\begin{tabular}{llll}
\toprule
Component & Type & Guarantee & Condition \\
\midrule
$\Pi$ & Def.~\ref{def:pipeline} & Type-safe composition & Morphism typing \\
$\Gamma$ & Def.~\ref{def:cgrag} & No $Q_0$ references admitted & By construction \\
$\mathcal{B}$ & Def.~\ref{def:billing} & Monotone cost accumulation & Linear coefficients \\
$\Phi_\theta$ & Def.~\ref{def:lcs} & Deterministic reconciliation & Fixed $\theta=4$ \\
\tvg{} & Thm.~\ref{thm:repair} & $P(\text{success}\mid K)=1-(1-p)^{K+1}$ & IID attempts \\
\texttt{concMap} & Thm.~\ref{thm:concurrent}(a) & Index-order output & Independent tasks \\
\texttt{concMap} & Thm.~\ref{thm:concurrent}(b) & $S(N)=1/(\alpha+(1-\alpha)/N)$ & Amdahl model \\
$\mathcal{B}$ & Prop.~\ref{prop:token} & $\mathbb{E}[W_K]=W_\infty(1-e^{-\lambda K})$ & Fixed $w$ per attempt \\
$\Pi$ & Prop.~\ref{prop:mdp} & MDP formulation exists & Markov property \\
\bottomrule
\end{tabular}
\end{table}
